\documentclass[11pt,letterpaper]{article}
\usepackage[T1]{fontenc}
\usepackage[utf8]{inputenc}
\usepackage{lmodern}
\usepackage[margin=1in,headheight=15pt]{geometry}
\usepackage{amsmath,amssymb,amsthm,mathtools}
\usepackage{booktabs,array,longtable,enumitem,microtype,xcolor}
\usepackage{listings,fancyhdr,hyperref}
\definecolor{ink}{HTML}{18354A}
\definecolor{muted}{HTML}{526773}
\hypersetup{colorlinks=true,linkcolor=ink,urlcolor=ink,citecolor=ink,pdftitle={Protected equations: a smaller iPRA proof kernel},pdfauthor={Research prepared for Vladimir Reshetnikov}}
\urlstyle{same}
\pagestyle{fancy}\fancyhf{}
\fancyhead[L]{\small\color{muted}Protected equations for iPRA}
\fancyhead[R]{\small\color{muted}11 November 2026}
\fancyfoot[C]{\thepage}
\setlength{\parskip}{4pt}\setlength{\parindent}{0pt}
\setlist{topsep=4pt,itemsep=2pt,parsep=0pt}
\newtheorem{theorem}{Theorem}[section]
\newtheorem{lemma}[theorem]{Lemma}
\newtheorem{proposition}[theorem]{Proposition}
\newtheorem{corollary}[theorem]{Corollary}
\theoremstyle{definition}\newtheorem{definition}[theorem]{Definition}
\theoremstyle{remark}\newtheorem{remark}[theorem]{Remark}
\newcommand{\D}{\mathsf D}\newcommand{\Q}{\mathsf Q}
\newcommand{\Hcal}{\mathcal H}\newcommand{\Kcal}{\mathcal K}
\newcommand{\Ncal}{\mathcal N}\newcommand{\Pcal}{\mathcal P}
\newcommand{\ipra}{\mathrm{iPRA}}
\newcommand{\eqv}{\mathrel{\dashv\vdash}}
\newcommand{\headI}{\mathsf I}
\newcommand{\file}[1]{{\urlstyle{tt}\nolinkurl{#1}}}
\lstset{basicstyle=\ttfamily\scriptsize,breaklines=true,columns=fullflexible,keepspaces=true,showstringspaces=false,frame=single,rulecolor=\color{muted},xleftmargin=3pt,xrightmargin=3pt}

\begin{document}
\begin{titlepage}
\vspace*{0.65in}
{\color{muted}\large CONCISE INTUITIONISTIC ARITHMETIC\par}
\vspace{0.2in}
{\color{ink}\Huge\bfseries Protected equations\par}
\vspace{0.12in}
{\color{ink}\LARGE A smaller iPRA proof kernel\par}
\vspace{0.35in}
{\large Construction, proof translations, and reproducible checker comparison\par}
\vspace{0.25in}
Prepared for Vladimir Reshetnikov\par
11 November 2026\par
\vspace{0.45in}
\begin{abstract}
The current ConciseIPRA collection contains two useful reductions: an
introduction rule whose choice is determined by the matrix, and an
elimination rule whose choice is determined by the substitution source.
Combining their programs alone does not establish a bi-interpretation:
the available proofs of completeness have incompatible equality
bootstraps. We resolve that issue by restricting public atomic formulas
to equations with an iterator at both roots. Ordinary equality is
represented by
\[
 W(s)=W(t),\qquad W(r):=I(0,r;u.0).
\]
All required proof transformations can then use equality sources with a
fixed iterator root. Such sources survive arbitrary numerical
substitution, including substitution of zero. This permits a
noncircular reconstruction of padding, numerical cases, equation
induction, and both interpretation composites.

Under the repository's canonical-statement measurement convention,
the three delivered complete checker profiles improve its recorded
coordinate minima from 259 to 258 significant tokens, from 283 to 281
AST nodes, and from 526 to 525 normalized bytes. These minima belong to
different programs. The improvement is small; no absolute minimality
or literature-wide novelty is claimed. A separate strict checker
validates the public statement grammar. The new equivalence proof is
written mathematics, supported by executable certificates; it has not
been formalized in Lean.
\end{abstract}
\vfill
\small Baseline frozen at commit\par
\file{1517d9d79ce33f631e0f82960c849427cbba23bf}\par
\href{https://github.com/VladimirReshetnikov/ProveIt/tree/1517d9d79ce33f631e0f82960c849427cbba23bf/Logic/PrimitiveRecursiveArithmetic/ConciseIPRA}{ProveIt / Logic / PrimitiveRecursiveArithmetic / ConciseIPRA}
\end{titlepage}
\tableofcontents
\newpage

\section{The result and its precise scope}

The proposed presentation is the \emph{protected core}. It uses the
same numerical terms as the current native scalar core, one public
formula binder, and a deterministic introduction schema. Its essential
new ingredient is a syntactic invariant on public equations, which
makes two previously separate reductions compatible.

\begin{theorem}[Main result]\label{thm:main}
The protected core is bi-interpretable and equiconsistent with full
first-order intuitionistic primitive recursive arithmetic. Both
numerical interpretation domains are whole. Interpreted equality is
actual equality, represented in the protected syntax by
\(W(x)=W(y)\). Both numerical composites are the identity up to this
equality. For every public formula, both formula composites are
provably equivalent to the original formula.
\end{theorem}

The comparison is against the pinned repository state, rather than an
earlier draft. Its README records the native calculus, its Lean
bi-interpretation, and the existing \file{source_dispatch} and
\file{guarded_core} packages~\cite{repo,source,guarded}. The latter two
packages supply the two reductions being combined. This report supplies
the protected-equation compatibility argument, a canonical language for
the combination, executable examples, and measured programs.

The source-size gain is deliberately not described as a large
breakthrough. It is a strict improvement under the existing comparison
contract. The strict public-grammar checker is larger than the earlier
guarded strict checker; the distinction is reported in
Section~\ref{sec:metrics}. A proof of global optimality would require
substantially more than comparison with the inspected implementations.

\section{The inherited arithmetic and native reference system}

There is one numerical sort. Terms are finite expressions
\[
 t ::= x\mid 0\mid S(t)\mid I(n,a;u.h).
\]
Only \(u\) in the update \(h\) is bound by the displayed iterator.
The counter \(n\), initializer \(a\), and passive parameters are not
bound by it. All substitutions below avoid capture. Put \(1=S0\).
The two computation schemas are
\begin{align}
 I(0,a;u.h)&=a,\tag{I0}\label{eq:i0}\\
 I(Sn,a;u.h)&=h[I(n,a;u.h)/u].\tag{IS}\label{eq:is}
\end{align}

The current Lean native system \(\Ncal\) has formulas
\(s=t,\Q_yB,\D_{x,y}B\), with ordinary readings
\[
 \Q_yB=B(0)\to\forall yB(y),\qquad
 \D_{x,y}B=\exists x\bigl(B(x,0)\to\forall yB(x,y)\bigr).
\]
Its eight constructors are hypothesis, I0, IS, equational transport,
QE, DE, witnessed case introduction of D, and equation-progress
introduction of Q. The repository proves transport in arbitrary
formula templates and witnessed D-progress as derived rules. It also
proves that these native rules present a theory bi-interpretable with
iPRA~\cite{native,bridge}.

For orientation, its ordinary arithmetic realization uses
intuitionistic equality logic, I0 and IS, numerical coverage,
\((0=1)\to\bot\), and contextual induction for literal equations.
The inherited arithmetic bridge derives predecessor, equality
reflection, quantifier-free induction, pairing, and a compiler for
every fixed primitive recursive definition. It proves the two
composites by internal uniqueness, not merely by standard-model
agreement. The relevant written proofs and selected Lean statements
are frozen in \file{reference/}.

We retain this arithmetic bridge. In particular, we do not assert that
positive computation equations alone, without a falsity bridge,
exclude a singleton ordinary structure.

\section{Complete specification of the protected core}\label{sec:theory}

\subsection{Public formulas and equality}

Call a term \emph{protected} when its outermost constructor is
\(I\). This is a syntactic condition, not a claim that its numerical
value is nonzero. Every protected term keeps its \(I\) root under
capture-free numerical substitution.

\begin{definition}[Public grammar]
The public formulas are generated by the following clauses.
\begin{enumerate}
\item \(p=q\) is a public formula if both \(p\) and \(q\) are protected.
\item \(\D_{x,y}B\) is public if \(B\) is a compound public formula.
\item An atomic matrix directly below D is allowed exactly in the form
\[
 B(x,y)=\bigl(I(y,0;u.h)=I(y,0;u.k)\bigr).
 \tag{Guard}\label{eq:guard}
\]
Thus \(\D_{x,y}B\) with this matrix is public. The updates may depend
on \(x,y\), the state \(u\), and passive parameters.
\end{enumerate}
\end{definition}

The bound variables \(x,y\) are distinct. There is no public Q
constructor. Define
\[
 W(t):=I(0,t;u.0),\qquad
 s\simeq t\quad:\!\!\equiv\quad W(s)=W(t),\qquad
 \bot_{\Pcal}:=(W(0)=W(1)).
\]
The variables bound in \(W\) are always fresh. The equality
\(\simeq\) is an abbreviation, not an additional primitive relation.
Since I0 proves \(W(t)=t\) in the ordinary realization, its reading
is actual equality throughout the entire numerical domain.

\subsection{Auxiliary judgments and the seven rule forms}

The proof calculus \(\Hcal\) also permits auxiliary equations with
arbitrary term endpoints and closure resources \(\Q_yB\). These are
needed, for example, to use the computation equation \(W(t)=t\)
when \(t\) itself is not protected. Auxiliary equations are not
additional public atomic predicates. A public theorem is a public
formula having a finite derivation in the stated calculus.

Retained contexts are finite reusable lists. They may include auxiliary
judgments in intermediate derivations. Q may be used with a public
matrix, and D formulas in \(\Hcal\) obey the public grammar. The
ordinary interpretation of all these auxiliary judgments is the
uniform conditional reading in Section~2.

Besides I0 and IS, available in every context, the rules are precisely
the following. \(F(v)\) is an allowed formula or auxiliary-equation
template, with the indicated instances well formed.
\[
 \frac{A\in\Gamma}{\Gamma\vdash A}\ (\mathrm{Hyp}),\qquad
 \frac{\Gamma\vdash p=q\quad\Gamma\vdash F(p)}
      {\Gamma\vdash F(q)}\ (\mathrm{Tr}_{I}),
 \quad p\text{ has an }I\text{ root}.
 \tag{H,Tr}\label{eq:tr}
\]
\[
 \frac{\Gamma\vdash\Q_yB(y)\quad\Gamma\vdash B(0)}
      {\Gamma\vdash B(t)}\ (\mathrm{QE}),
\]
\[
 \frac{\Gamma\vdash\D_{x,y}B(x,y)\quad
       \Gamma,\Q_yB(b,y)\vdash A}
      {\Gamma\vdash A}\ (\mathrm{DE}).
\]
In DE, \(b\) is fresh for the retained context, conclusion, and fixed
parameters. The sole D-introduction schema is
\[
 \frac{\Gamma,B(t,\eta_B(z))\vdash B(t,Sz)}
      {\Gamma\vdash\D_{x,y}B(x,y)},\qquad
 \eta_B(z)=
 \begin{cases}
 z,&B\text{ is an equation},\\
 0,&B\text{ is compound}.
 \end{cases}
 \tag{DI}\label{eq:di}
\]
Here the conclusion must be public, \(z\) is fresh for \(\Gamma,t\)
and the fixed parameters, and the displayed local resource is
discharged. At an equation matrix the guard requirement
\eqref{eq:guard} applies. The certificate does not choose between
cases and progress; the matrix root determines the premise.

Counting the two computation schemas separately gives seven forms:
Hyp, I0, IS, Tr, QE, DE, and DI. There are no primitive reflexivity,
successor-injectivity, separation, coverage, conventional connective,
or induction axioms. Structural rules are admissible operations on
finite derivations, not additional certificate instructions.

\subsection{The executable extension \texorpdfstring{\(\Kcal\)}{K}}

The compact programs permit more administrative syntax. In particular,
their transport may use any source other than the literal zero, and
their raw D formulas need not pass the public guard test. Introduction
is still determined by the matrix root. This extension is denoted
\(\Kcal\); every \(\Hcal\) inference has an encoding in it.

Every such raw inference is sound under the \emph{same ordinary
conditional reading}: equation-progress introduction is sound by
equation induction under its conditional base; compound-case
introduction is sound by numerical coverage. No progressive alternative
meaning of D is necessary for our soundness argument. The restrictions
on public formulas are needed for reverse translation and formula
composites. We do not claim completeness or a bi-interpretation for
every unrestricted raw statement.

\section{The safe equality bootstrap}\label{sec:safe}

The key invariant is stronger than ``the source is not zero'': every
transport source used in the completeness proof has an \(I\) root.
This makes the proof transformations stable under substitution.

\begin{lemma}[Structural operations]\label{lem:struct}
Weakening, exchange, contraction, fresh renaming, numerical substitution,
and ordinary-assumption cut are admissible in \(\Hcal\).
\end{lemma}
\begin{proof}
Induct on derivations, renaming eigenparameters before substitution or
cut. Public guards retain their shape because their bound counter is
not substituted by free-parameter substitution. An \(I\)-headed
transport source retains its head. All other rule side conditions are
the usual explicit freshness conditions. Cut replaces a hypothesis use
by its supplied derivation. This proof would fail for a rule merely
requiring a variable source to be syntactically nonzero; that weaker
invariant is not being used here.
\end{proof}

\begin{lemma}[Safe equality operations]\label{lem:safe}
The following operations have derivations whose transport sources all
have \(I\) roots.
\begin{enumerate}
\item Reflexivity \(t=t\) for any term \(t\).
\item From \(p=q\), symmetry \(q=p\), when \(p\) is protected.
\item From \(r=p\) and \(p=q\), transitivity \(r=q\), when \(p\)
is protected.
\item From \(p=q\), the equation \(f(p)=f(q)\), when \(p\) is
protected and \(f\) is a fixed numerical template.
\end{enumerate}
Consequently ordinary equality algebra is available between protected
endpoints.
\end{lemma}
\begin{proof}
Put \(a=W(t)\). I0 gives \(a=t\). Use this equation as both
premises of transport with template \(v=t\); the result is \(t=t\).
For symmetry transport \(p=p\) along \(p=q\) in the template
\(v=p\). For transitivity use the template \(r=v\). For congruence
transport reflexivity of \(f(p)\) in the template \(f(p)=f(v)\).
The last operation may produce a nonprotected endpoint; such an
endpoint is never silently used as a transport source in later
arguments.
\end{proof}

\begin{lemma}[Joining at zero]\label{lem:zerojoin}
For protected \(p,q\), from \(p=0\) and \(q=0\) one can derive
\(p=q\) using only \(I\)-headed transport sources.
\end{lemma}
\begin{proof}
Let \(w=W(0)\). From \(p=0\), transport reflexivity
\(W(p)=W(p)\) on the left to obtain \(w=W(p)\).
Its safe symmetry gives \(W(p)=w\). Now use the computation
\(W(p)=p\) to obtain \(p=w\). Similarly obtain \(q=w\).
All the endpoints \(p,q,w\) are protected, so safe symmetry and
transitivity yield \(p=q\). At no point is \(0=p\) used as a
transport major.
\end{proof}

\begin{lemma}[Lifting an arbitrary equality]\label{lem:lift}
For arbitrary terms \(s,t\) and a numerical template \(a(v)\),
\[
 W(s)=W(t)\ \vdash_{\Hcal}\ W(a(s))=W(a(t)).
\]
Consequently the native equational transport rule translates under
\(s=t\mapsto W(s)=W(t)\).
\end{lemma}
\begin{proof}
Congruence on the given protected equation gives
\(W(a(W(s)))=W(a(W(t)))\). The two computation equations
\(W(s)=s\) and \(W(t)=t\), whose sources are protected, give
\[
 W(a(W(s)))=W(a(s)),\qquad
 W(a(W(t)))=W(a(t)).
\]
Every endpoint in these three displayed equations has a fixed \(W\)
root. They can therefore be reversed and composed safely, proving
the claim. Apply this construction to both numerical sides of an old
equation template, and compose with the translated minor premise.
The argument still works when \(s\) is literally zero.
\end{proof}

This is the compatibility step absent from a bare composition of the
two old programs. It replaces a demand for unrestricted zero-source
symmetry by a translation of all public equalities into protected
ones.

\section{Padding, numerical cases, and equation induction}\label{sec:induction}

Set \(C(v,a,b)=I(v,a;u.b)\), with the state variable unused. For a
protected equation \(E=(p=q)\), define
\[
 G_vE=(C(v,0,p)=C(v,0,q)),\qquad
 L(E)=\D_{a,v}G_vE,
\]
where \(a,v\) are fresh for \(E\). Let \(P(E)=L(E)\) for an
equation, and \(P(B)=B\) for a compound public formula. Thus
\(P(B)\) is always compound and public.

\begin{lemma}[Guard endpoints and padding]\label{lem:pad}
For every protected equation \(E\),
\[
 \vdash G_0E,\qquad G_{Sz}E\eqv E,\qquad L(E)\eqv E.
\]
These interderivabilities are uniform under additional retained context.
\end{lemma}
\begin{proof}
The two sides of \(G_0E\) compute to zero. They are protected,
so Lemma~\ref{lem:zerojoin} joins them. At a successor, the two sides
compute to \(p\) and \(q\). Both are protected, so safe equality
algebra gives the endpoint equivalence in both directions.
Under \(E\), atomic DI proves \(L(E)\): its successor goal follows
from \(E\), independently of its local induction assumption.
Conversely, DE on \(L(E)\) supplies \(\Q_vG_vE\). Feed its proven
zero instance and use QE at \(1\), then remove the successor guard.
\end{proof}

For protected equations \(U=(p=q)\), \(V=(r=s)\), put
\[
 M_v(U,V)=(C(v,p,r)=C(v,q,s)),\qquad
 U\Rightarrow_g V:=\D_{a,v}L(M_v(U,V)).
\]
The outer matrix is compound. Endpoints and padding show that this
macro has implication introduction and elimination between protected
equations, using compound-case DI and DE followed by QE at \(1\).

\begin{lemma}[Numerical cases]\label{lem:cases}
For every public formula \(A(y)\),
\[
 \frac{\Gamma\vdash A(0)\qquad\Gamma\vdash A(Sz)}
      {\Gamma\vdash A(t)}
\]
is admissible when \(z\) is fresh for the retained context and fixed
parameters.
\end{lemma}
\begin{proof}
If \(A\) is compound, introduce \(\D_{a,y}A(y)\) with unused
\(a\) by cases, ignoring its local base assumption. Eliminate it,
supply the given base, and use QE at \(t\). For an equation apply
this compound case construction to \(L(A(y))\), padding both
premises and unpadding the conclusion.
\end{proof}

\begin{theorem}[Contextual induction for protected equations]\label{thm:ei}
For every equation \(E(y)\) whose two sides have fixed \(I\) roots,
\[
 \frac{\Gamma\vdash E(0)\qquad\Gamma,E(z)\vdash E(Sz)}
      {\Gamma\vdash E(t)}
 \tag{PEI}\label{eq:pei}
\]
is admissible, with the usual eigenvariable condition.
\end{theorem}
\begin{proof}
Put \(J(y)=G_yE(y)\). We need its progress step without retaining
a case-variable-dependent assumption during numerical cases.
Use the compound formula
\(\Phi(y)=J(y)\Rightarrow_g E(Sy)\).

For \(y=0\), ignore the implication's local assumption, take the
given \(E(0)\), and specialize the given step at zero to obtain
\(E(1)\). For \(y=Sw\), the assumption \(J(Sw)\) gives
\(E(Sw)\); the step specialized at \(Sw\) gives \(E(SSw)\).
Both branch derivations have retained context \(\Gamma\), independent
of the case variable. Lemma~\ref{lem:cases} yields \(\Phi(z)\).
Under \(J(z)\), implication elimination therefore yields \(E(Sz)\),
and the successor guard gives \(J(Sz)\).

Atomic DI now proves \(\D_{a,y}J(y)\) with unused \(a\).
DE and the theorem \(J(0)\) supply \(J(Sw)\), hence \(E(Sw)\),
for fresh \(w\). Numerical cases combine these successor instances
with the original \(E(0)\), yielding \(E(t)\).

The substitutions into the supplied step are legitimate by
Lemma~\ref{lem:struct}: its derivation is in \(\Hcal\), whose
transport sources retain an \(I\) root. This is the point at which a
naive appeal to arbitrary nonzero-source derivations would be circular.
\end{proof}

\section{Explicit translations and the two round trips}\label{sec:translations}

\subsection{Translation from the formalized native calculus}

Define \(T\) on finite native formulas, with all new binders fresh:
\begin{align*}
 T(s=t)&=W(s)=W(t),\\
 T(\Q_y B)&=\D_{a,y}P(TB)\quad(a\text{ unused}),\\
 T(\D_{x,y}B)&=\D_{x,y}P(TB).
\end{align*}
Numerical terms are unchanged. This translation is effective, commutes
with capture-free numerical substitution up to renaming, and always
produces public formulas. Translate contexts pointwise.

Write \(\operatorname{Cl}_y B=\D_{a,y}P(B)\) for the translated
closure. From \(\operatorname{Cl}_yB\) and \(B(0)\), DE, padding,
QE, and unpadding derive \(B(t)\). From
\(\Gamma,B(0)\vdash B(Sz)\), compound-case DI proves
\(\operatorname{Cl}_yB\). These statements use the endpoint
interderivabilities of Lemma~\ref{lem:pad}.

\begin{theorem}[Native proof translation]\label{thm:translate}
\[
 \Gamma\vdash_{\Ncal}A\quad\Longrightarrow\quad
 T\Gamma\vdash_{\Hcal}TA.
\]
The transformation is effective on finite derivations and uses only
\(I\)-headed transport sources.
\end{theorem}
\begin{proof}
Induct on the native derivation. Hypothesis use is immediate. Each
computation equation has an \(I\)-headed source, so safe congruence
through \(W\) gives its translated computation. Native equational
transport is Lemma~\ref{lem:lift}, applied to both sides. QE uses
the translated closure elimination just described.

For native case DI, unpad its translated local base, apply the
translated premise, and pad the successor result. The target matrix
\(P(TB)\) is compound, so the root-determined rule uses cases.

For native Q-progress on \(E\), its translated step is a step for
the protected equation \(TE\). Under the translated base, use
Theorem~\ref{thm:ei} at \(Sz\). Pad the result and use compound
case introduction to prove \(\operatorname{Cl}_yTE\).

For native DE, its translated major opens the raw resource
\(\Q_yP(TB(b,y))\). This resource supplies its corresponding public
closure \(\operatorname{Cl}_yTB(b,y)\): introduce that closure by
cases, and use the retained raw Q resource at the successor. Cut this
public closure into the recursively translated minor derivation, then
close DE. Freshening all auxiliary binders preserves the original
eigenvariable conditions.
\end{proof}

\subsection{Literal inclusion and the canonical round trip}

Let \(J\) send a public formula to the identical native Eq/D formula.
Every \(\Hcal\) inference is valid in \(\Ncal\): unrestricted
formula transport and witnessed equation-progress are derived native
rules. This provides the reverse proof interpretation.

\begin{lemma}[Public round trip]\label{lem:roundtrip}
For every public formula \(A\),
\[
             A\eqv_{\Hcal}T(JA).
\]
\end{lemma}
\begin{proof}
Induct on the number of D binders. For an equation \(p=q\), both
endpoints are protected. The computations \(W(p)=p\) and
\(W(q)=q\) therefore have protected endpoints, and safe equality
algebra proves
\(p=q\eqv W(p)=W(q)\).

For a D formula with compound matrix, use compound binder congruence:
eliminate the source at a fresh witness; introduce the target by
cases; convert its base backwards, use the opened closure at the
successor, and convert the result forwards. Matrix interderivability
is given by the external induction hypothesis. All conversions use
safe derivations.

It remains to treat \(\D_{x,y}E\) with an atomic guard matrix.
Its zero instance \(E(x,0)\) is a theorem, uniformly in \(x\), by
Lemma~\ref{lem:zerojoin}. Its translation is
\(\D_{x,y}L(TE)\).
Forwards, DE opens \(\Q_yE(b,y)\). Supply its proven zero instance,
obtain each successor instance, translate the protected equation by
the atomic equivalence, and pad. Compound-case DI at witness \(b\)
then proves \(\D_{x,y}L(TE)\).
Backwards, DE opens \(\Q_yL(TE(b,y))\). Pad the translated zero
theorem and use QE at \(Sz\). Unpadding and the atomic equivalence
give \(E(b,Sz)\). Atomic-progress DI, ignoring its local induction
assumption, proves \(\D_{x,y}E\). Close DE in both directions.
This case does not assume a nonexistent atomic case-introduction rule.
\end{proof}

The other native formula composite follows by external induction:
native equality proves \(W(t)=t\), native padding is equivalent to
its matrix, and the unused-witness D represents Q. Thus
\(JT(A)\eqv_{\Ncal}A\) for every native formula, including those
with Q constructors.

\begin{corollary}[Equality of canonical consequence relations]\label{cor:canonical}
For finite public contexts \(\Gamma\) and public \(A\),
\[
 \Gamma\vdash_{\Hcal}A
 \quad\Longleftrightarrow\quad
 J\Gamma\vdash_{\Ncal}JA.
\]
For empty external context and closed canonical goals, the same
equivalence holds for finite accepting compact certificates under the
decoding contract in Section~\ref{sec:checker}.
\end{corollary}
\begin{proof}
Left to right is literal native soundness. For the converse, translate
the native proof by Theorem~\ref{thm:translate}, then use the public
round trips on the retained assumptions and conclusion, with cut.
The compact extension is native-sound, and every \(\Hcal\) proof
has an encoding, so its canonical consequences lie between these
equal consequence relations.
\end{proof}

\section{Full intuitionistic logic and bi-interpretability with iPRA}

There is no restriction to the equational or quantifier-free fragment
of iPRA. To specify the logical interface explicitly, let
\(\operatorname{enc}\) be the existing native encoding of ordinary
formulas, and define
\[
             \operatorname{enc}_{\Pcal}(F):=T(\operatorname{enc}(F)).
\]
This is a finite structural macro expansion, not an appeal to a
provability oracle. The native encoding constructs implication,
conjunction, disjunction, existence, and universality from its selector
and gate operations~\cite{native}. For example, its macros are
\begin{align*}
 A\to B&:=\Q_y\operatorname{sel}(y,A,B),\\
 A\wedge B&:=\Q_y\operatorname{gate}
       (y,\operatorname{sel}(\operatorname{par}(y),A,B)),\\
 \exists xA&:=\D_{x,y}\operatorname{gate}(y,A(x)),\\
 A\vee B&:=\exists v\operatorname{sel}(v,A,B),\\
 \forall xA&:=A(0)\wedge\Q_yA(y).
\end{align*}
Here the formula selector and gate are defined by recursion on finite
formula syntax, using the numerical selector \(C\), binder congruence,
and padding. Applying \(T\) removes the public Q constructors and
protects all atoms. All introduction and elimination rules transfer by
Theorem~\ref{thm:translate}. Thus the displayed complete native
macro implementation in \file{reference/Lean/NativeLogic.lean},
composed with our explicit \(T\), specifies all ordinary logical
operations without adding axioms or trusted checker code.

For numerical terms, the forward map to iPRA uses the inherited
primitive-recursive implementation of each fixed iterator. The reverse
map uses its scalar compiler for each fixed primitive-recursive
definition. No universal evaluator over arbitrary function codes is
postulated. The inherited round-trip equalities translate by \(T\)
to protected equalities, so they hold internally in the new theory.

The domain predicate in both directions is the whole domain. The
numerical identity graph is \(x=y\) in the ordinary realization and
\(W(x)=W(y)\) in the protected presentation. Since I0 identifies
\(W(t)\) with \(t\) ordinarily, there is no quotient, retract,
standard-numeral restriction, or hidden change in equality.
The public round trip and the inherited native/iPRA round trips prove
both full formula composites. Finally,
\[
 \Pcal\vdash W(0)=W(1)
 \quad\Longleftrightarrow\quad
 \Ncal\vdash0=1
 \quad\Longleftrightarrow\quad
 \ipra\vdash\bot.
\]
This proves Theorem~\ref{thm:main}. It is a metatheoretic
equiconsistency statement, not a claim that either theory proves its
own consistency.

\section{Why the protected language matters}

The words ``not literally zero'' alone do not supply the structural
invariant needed in the proof. A variable source is permitted by that
test, but substituting zero can turn its use into a forbidden rule
instance. Moreover, guarding an arbitrary equation can require
inserting wrappers around a zero endpoint; unrestricted equality can
do this, whereas restricted forward transport cannot simply assume it.
Using one old completeness argument inside the other without checking
these dependencies would be circular.

Protected atoms resolve both issues. All equality information that
must later act as a major is represented at stable iterator roots.
Zero computations are joined through \(W(0)\), and arbitrary native
equalities are moved through the sandwich in
Lemma~\ref{lem:lift}. Numerical substitution is first proved in
\(\Hcal\), before it is used in the induction bootstrap.

There is a related distinction in the upstream README. The statement
that substitution is not \emph{admissible} in its already-equivalent
\(\mathrm{ZPrf}\) system is too strong. Its equivalence theorem
\file{zprf_iff_nprf} transports native substitution to ZPrf. What
fails is direct substitution into the same derivation tree without
repairing rule instances. This correction does not solve the guarded
composition by itself, since equivalence of that new composition has
not yet been established at that stage. Our stable-head proof avoids
using such a corollary prematurely.

\section{Executable certificate contract}\label{sec:checker}

\subsection{Trusted programs and input}

Each compact profile is a complete standalone Python module importing
only the standard JSON parser. It reads a JSON array from standard
input. The last element is the closed goal; preceding instruction
packets are consumed backwards. Success is normal process exit with
exactly \file{ACCEPT} followed by a newline. An exception, nonzero
exit, timeout, or resource exhaustion is not acceptance. Unused prefix
data is inert, as in the inherited protocol. No producer, normalizer,
theorem database, or Python text supplied by the certificate is executed.

The arity encoding is inherited. At numerical positions, unary arrays
mean successor and ternary arrays mean iterator. At formula positions,
unary arrays mean D, binary arrays mean equality, and ternary arrays
mean Q. A canonical Q has three copies of its body. The opener adjusts
depth by \(5-\text{array arity}\). Protected slots have width two;
numerical root depth is three and formula root depth zero. The full
sorted interpretation, including zero and truth aliases, is supplied
in \file{reference/inherited_arity_checker.tex}.

The public grammar adds two explicit requirements to this convention:
both numerical children of every public equation have ternary iterator
roots; and an atomic D matrix has both sides beginning with
\file{[5, 0]}, the inherited encoding of the displayed guard. The
strict checker enforces the closed canonical goal grammar before
running the same kernel. It does not reject sound administrative
intermediate resources merely because they are not public formulas.

\subsection{The two combined changes}

The source-dispatch matcher is
\begin{lstlisting}
m == ([s, u] if s else [t] * 3)
\end{lstlisting}
so the source determines whether the major is equality or Q.
The root-determined introduction uses
\begin{lstlisting}
(m or not len(a) % 2) and k
\end{lstlisting}
so an atomic matrix takes its progress instance and a compound matrix
takes its zero instance. The latter removes the old comparison with a
certificate-supplied choice. The complete byte profile appears in
Appendix~\ref{app:code}.

Canonical \(\Hcal\) transport sources have a ternary iterator root
and are therefore truthy in this encoding. The source matcher accepts
them. Atomic progress and compound case introduction match the new
root test. The inherited computation normalization supplies the
restricted computation packets while preserving their conclusions;
public guards need not be changed. A producer and concrete safe
equality normalization are delivered, but are outside the trusted
program.

\subsection{Soundness and the scope of completeness}

Soundness of a finite accepting trace uses the inherited opening,
scope, and sorted-alias lemmas. The source-selected equality and Q
major are still verified. The new introduction chooses one of the two
already-sound Choice-DI rules. Thus every accepted raw proof has a
native-provable ordinary reading.

For a canonical public goal, Corollary~\ref{cor:canonical} turns that
sound reading back into a protected proof. Conversely every finite
\(\Hcal\) proof can be encoded with fresh integer parameters and
the inherited computation expansions. This proves equality of
canonical theorems in the ideal unbounded implementation. Physical
Python stack and memory limits can reject large valid proofs; no
completeness claim over a fixed finite resource bound is made.

The compact profiles intentionally do not validate the exact public
goal grammar. Their record is measured under the same
canonical-statement convention used by the source packages. The
strict profile supplies the additional public-boundary check and is
counted in its entirety. We do not omit that validator while claiming
an exact-public-grammar program of 525 bytes.

\section{Reproducible size comparison}\label{sec:metrics}

The measurements use CPython 3.12.14, the version used by the current
repository. Normalized source is \file{ast.unparse} output with a
final newline. Comments and non-executable docstrings do not contribute
to significant-token or executable-AST counts. Every selected compact
program includes JSON input, scope traversal, proof checking, and
output. Gzip uses level 9 and \file{mtime=0}; it is an informational
coordinate, not an optimization objective.

\begin{center}\small
\begin{tabular}{llrrrr}
\toprule
Profile & Program & Tokens & AST nodes & Bytes & Gzip-9\\
\midrule
Token & prior source-dispatch & 259 & 289 & 572 & 339\\
      & protected core & \textbf{258} & 287 & 571 & 339\\
AST   & prior source-dispatch & 263 & 283 & 607 & 348\\
      & protected core & 262 & \textbf{281} & 606 & 348\\
Byte  & prior source-dispatch & 266 & 305 & 526 & 330\\
      & protected core & 265 & 303 & \textbf{525} & 331\\
\bottomrule
\end{tabular}
\end{center}

Thus the coordinate envelope decreases from \((259,283,526)\) to
\((258,281,525)\). It is not the size triple of one program.
The token and AST profiles improve their corresponding old programs
in all three principal coordinates; the byte profile does too.
Compression of the byte profile increases by one byte, which is
reported rather than hidden.

The prior guarded profiles have minima 260 tokens, 286 AST nodes,
and 527 normalized bytes. Their combination with source dispatch
therefore improves that frontier as well. This is a comparison with
the recorded current implementations, not a lower-bound theorem for
every possible proof checker.

The strict protected checker has 505 significant tokens, 582 AST
nodes, 1,251 normalized bytes, and 582 gzip-9 bytes. The previous
guarded strict checker has 1,208 normalized bytes. There is therefore
\emph{no strict-parser size improvement} in this comparison. The new
global protected-atom check has a real cost. All raw and normalized
measurements, frozen source hashes, and generated gzip files are
included in the evidence directory.

\section{Validation and research status}\label{sec:validation}

The reproduction suite passed all 608 executions. It checks 5 candidate programs and 6 frozen comparison programs, using 11 named canonical examples and 22 negative fixture families. Candidate executions cover both ordinary Python and Python with optimization enabled. The examples include both wrapped computation schemas, protected reflexivity, compound introduction and elimination, both padding directions, a 564-instruction guarded induction certificate, and a guarded implication expressing explosion from the protected contradiction. All generated named derivations pass the guard and source audits.


The induction example is not a vacuous use of its induction
hypothesis: the supplied audit records an actual use and verifies
that the safe normalization introduces no literal-zero transport
source. The certificates preserve the guarded matrices. The
contradiction example uses the public contradiction
\(W(0)=W(1)\), rather than relying on an unprotected equation
\(0=1\) as an interchangeable raw assumption.

These executable results verify the delivered programs and examples.
They do not prove the universal interpretation theorem by finite
testing. The theorem is supported by the explicit structural proof
above and an independent written review supplied in the archive.
The inherited native/iPRA bridge is formalized in the cited repository;
its Lean build was not rerun here. The protected-core translation and
the small raw Python kernel have not been verified in Lean or Rocq.

\section{Using the archive}

From the archive root, run
\begin{lstlisting}
python3 reproduce.py
pdflatex -interaction=nonstopmode -halt-on-error report.tex
pdflatex -interaction=nonstopmode -halt-on-error report.tex
\end{lstlisting}
The Python reproduction uses the standard library. It regenerates
certificates, audits them, executes positive and negative controls,
and reports size measurements. A TeX installation with the packages
listed in the preamble is needed only to rebuild the PDF.

\begin{center}\small
\begin{tabular}{ll}
\toprule
File & Purpose\\
\midrule
\file{code/checker_bytes.py} & Smallest byte profile\\
\file{code/checker_tokens.py} & Smallest token profile\\
\file{code/checker_ast.py} & Smallest AST profile\\
\file{code/readable_checker.py} & Expanded control flow for review\\
\file{code/strict_checker.py} & Public-goal validation plus kernel\\
\bottomrule
\end{tabular}
\end{center}
Local-scope and global-scope certificate files are labeled separately;
they are not claimed interchangeable.

The reproducibility manifest identifies the precise original files
and Git blobs used in the comparison. The inherited references retain
their provenance. The novel contribution is the protected language,
stable-source completeness proof, compatible combined kernels, and
their validation. No repository files were changed remotely.

\appendix
\section{The complete byte-profile checker}\label{app:code}

\lstinputlisting{code/checker_bytes.py}

\section{A compact mathematical inventory}

\begin{center}\small
\begin{tabular}{p{0.25\linewidth}p{0.66\linewidth}}
\toprule
Component & Specification\\
\midrule
Numerical terms & \(x,0,S(t),I(n,a;u.h)\); one numerical sort.\\
Public atoms & \(p=q\), both endpoints with a fixed iterator root.\\
Public binder & \(\D_{x,y}B\); an atomic matrix must have the zero guard.\\
Equality macro & \(s\simeq t\) abbreviates \(W(s)=W(t)\).\\
Contradiction & \(W(0)=W(S0)\).\\
Auxiliary resources & Arbitrary equations and conditional closures Q.\\
Axioms & The two computation schemas I0 and IS.\\
Rules & Hyp, stable-source Tr, QE, DE, and deterministic DI.\\
Derived structure & Numerical substitution and cut; no structural kernel rules.\\
Derived arithmetic & Protected equation induction, then the translated native arithmetic bridge.\\
Domain composite & Whole numerical domain; identity graph represented by \(\simeq\).\\
Formal status & New written proof and executable regression evidence; inherited Lean bridge.\\
\bottomrule
\end{tabular}
\end{center}

\begin{thebibliography}{9}
\bibitem{repo} ProveIt contributors.
\emph{Concise presentations of intuitionistic primitive recursive arithmetic}.
README at commit \file{1517d9d79ce33f631e0f82960c849427cbba23bf}.
\href{https://github.com/VladimirReshetnikov/ProveIt/tree/1517d9d79ce33f631e0f82960c849427cbba23bf/Logic/PrimitiveRecursiveArithmetic/ConciseIPRA}{Pinned repository version}.

\bibitem{native} ProveIt contributors.
\emph{Native.lean, NativeStruct.lean, NativeEq.lean, NativeLogic.lean}.
Same pinned commit, \file{Lean/ConciseIPRA/}. Selected exact sources
are frozen in the accompanying \file{reference/Lean/} directory.

\bibitem{bridge} ProveIt contributors.
\emph{NativeBiInterpretation.lean} and the inherited ordinary arithmetic
bridge. Same pinned commit. The former states the faithful proof
translations, both full formula round trips, and equiconsistency;
the latter is reproduced with the predecessor, equality-reflection,
and primitive-recursion arguments in \file{reference/}.

\bibitem{source} ProveIt contributors.
\emph{Source-dispatch research package} and \emph{NativeZeroSource.lean}.
Same pinned commit, \file{Research/source_dispatch/} and
\file{Lean/ConciseIPRA/NativeZeroSource.lean}.
\href{https://github.com/VladimirReshetnikov/ProveIt/tree/1517d9d79ce33f631e0f82960c849427cbba23bf/Logic/PrimitiveRecursiveArithmetic/ConciseIPRA/Research/source_dispatch}{Pinned source-dispatch package}.

\bibitem{guarded} ProveIt contributors.
\emph{A guarded core with a deterministic introduction rule}.
Same pinned commit, \file{Research/guarded_core/}. The guarded
bootstrap is inherited; the restriction to protected atoms and the
stable-source compatibility argument are the new work here.
\href{https://github.com/VladimirReshetnikov/ProveIt/tree/1517d9d79ce33f631e0f82960c849427cbba23bf/Logic/PrimitiveRecursiveArithmetic/ConciseIPRA/Research/guarded_core}{Pinned guarded-core package}.
\end{thebibliography}
\end{document}
